\apendice{Formulário Amostragem de Experiência}
\label{ap:forms-esm}
\newcolumntype{L}{>{\raggedright\arraybackslash}p{4.8cm}}
\newcolumntype{C}{>{\centering\arraybackslash}p{0.7cm}}
 
% =====================================================
% COMANDO PARA GERAR UMA TABELA DE ESCALA LIKERT (1 A 5)
% Uso: \escalalikert{legenda}{ancora esquerda}{ancora direita}
% =====================================================
\newcommand{\escalalikert}[3]{%
\begin{table}[H]
\centering
\caption{#1}
\begin{tabular}{L C C C C C L}
\toprule
\textbf{#2} & 1 & 2 & 3 & 4 & 5 & \textbf{#3} \\
\bottomrule
\end{tabular}
\legend{Fonte: Dados da pesquisa (2026).}
\end{table}
}

\section*{Apresentação}
Este questionário avalia a percepção do(a) respondente sobre usabilidade, imersão e desenvolvimento do Pensamento Computacional durante o uso da plataforma. As questões estão organizadas em três dimensões: \textit{Imersão e Foco}, \textit{Uso do Sistema (Usabilidade e Interface)} e \textit{Pensamento Computacional (PC)}. Em cada item, a resposta deve ser registrada em uma escala do tipo \mbox{\textit{Likert}}, variando de 1 a 5, conforme as âncoras apresentadas nos extremos da escala.
 
\section{Dimensão 1: Imersão e Foco}
 
\subsection{Concentração Profunda (Esforço de Atenção)}
\textit{(Resposta obrigatória)}
 
O quanto você se sentiu focado na resolução do problema enquanto jogava?
 
\escalalikert{Escala de concentração profunda}{Totalmente distraído (pensando em outras coisas fora do software)}{Totalmente absorvido (foco 100\% direcionado ao desafio)}
 
\subsection{Distorção Temporal (Percepção do Tempo)}
\textit{(Resposta obrigatória)}
 
Durante a realização, como foi a sua percepção sobre a passagem do tempo?
 
\escalalikert{Escala de distorção temporal}{O tempo passou muito devagar (fase arrastada e cansativa)}{O tempo voou (nem percebi o tempo passando de tão envolvido)}
 
\section{Dimensão 2: Uso do Sistema (Usabilidade e Interface)}
 
\subsection{Clareza do Sistema (Entendimento das Regras)}
\textit{(Resposta obrigatória)}
 
Ficou claro o que era para ser feito para avançar ou vencer as fases?
 
\escalalikert{Escala de clareza do sistema}{Totalmente confuso (não entendi os comandos nem o objetivo do jogo)}{Totalmente claro (entendi perfeitamente as mecânicas e os botões)}
 
\subsection{Facilidade de Uso (Esforço de Interação)}
\textit{(Resposta obrigatória)}
 
Como você avalia a facilidade de interagir com os elementos da tela e responder às questões?
 
\escalalikert{Escala de facilidade de uso}{Muito difícil/Travado (a interface atrapalhou minha linha de raciocínio)}{Muito fácil/Fluido (consegui executar os comandos sem pensar na interface)}
 
\subsection{Qualidade do Feedback (Resposta do Software)}
\textit{(Resposta obrigatória)}
 
Quando você cometeu um erro ou acertou, o sistema informou isso de maneira compreensível e útil?
 
\escalalikert{Escala de qualidade do feedback}{Não, fiquei perdido sem saber se o sistema aceitou ou não minha resposta}{Sim! O sistema me deu um retorno visual/textual imediato e inteligente}
 
\section{Dimensão 3: Pensamento Computacional (PC)}
Pense em como você resolveu este desafio no software em comparação a como resolve um exercício comum no caderno:
 
\subsection{Decomposição (Dividir para resolver)}
\textit{(Resposta obrigatória)}
 
Para resolver este desafio no software, você precisou quebrar o problema matemático em etapas ou partes menores?
 
\escalalikert{Escala de decomposição}{Não, tentei resolver tudo de uma vez só}{Sim! Dividir o problema em pequenos passos facilitou encontrar a resposta}
 
\subsection{Reconhecimento de Padrões (Identificar semelhanças)}
\textit{(Resposta obrigatória)}
 
Você percebeu alguma lógica ou estrutura na fase que já tinha visto em desafios anteriores do jogo?
 
\escalalikert{Escala de reconhecimento de padrões}{Não, cada fase parece um problema completamente novo e sem relação}{Sim! Identificar o padrão me ajudou a resolver muito mais rápido do que em uma aula comum}
 
\subsection{Abstração (Focar no que importa)}
\textit{(Resposta obrigatória)}
 
Ao analisar o problema na tela, você conseguiu focar apenas nas variáveis e dados matemáticos essenciais?
 
\escalalikert{Escala de abstração}{Não, fiquei confuso com os elementos do jogo e me perdi na matemática}{Sim! O formato visual me ajudou a isolar os dados importantes melhor do que ler um enunciado de texto tradicional}
 
\subsection{Algoritmos e Depuração (Criar e corrigir sequências)}
\textit{(Resposta obrigatória)}
 
Para vencer a fase, você teve que planejar uma sequência lógica de comandos/passos e conseguiu corrigir seus próprios erros?
 
\escalalikert{Escala de algoritmos e depuração}{Não, resolvi na base do chute ou tentativa e erro sem planejamento}{Sim! Desenvolvi um passo a passo lógico e consegui corrigir meus erros de forma autônoma (depurar)}
 
\subsection{Retenção e Utilidade Percebida}
\textit{(Resposta obrigatória)}
 
Você sente que a lógica usada neste desafio ajudará você a lembrar das regras matemáticas na hora da prova mais do que se tivesse apenas copiado do quadro?
 
\escalalikert{Escala de retenção e utilidade percebida}{Não, acho que vou esquecer mais rápido do que o método tradicional}{Sim! Praticar interagindo fixou o assunto de forma muito mais profunda}
 
